% =====================================================
% 题型：三角函数图像性质选择题 (q_trig_graph_properties.tex)
% =====================================================
% 难度：★★ (中档)
% =====================================================
\newcommand{\qTrigGraphPropertiesStem}{已知函数 \(f(x)=2\sin\left(2x+\dfrac{\pi}{6}\right)\)，则下列说法正确的是 \hfill （\quad）}

\newcommand{\qTrigGraphPropertiesOptions}{%
  \mcoptions[1]{函数的最小正周期为 \(\pi\)}{函数图像关于直线 \(x=\dfrac{\pi}{6}\) 对称}{函数在区间 \(\left[-\dfrac{\pi}{3},\dfrac{\pi}{6}\right]\) 上单调递增}{函数图像向右平移 \(\dfrac{\pi}{12}\) 个单位可得到 \(y=2\sin2x\) 的图像}%
}

\newcommand{\qTrigGraphPropertiesFigure}[1][1.0]{%
  \begin{tikzpicture}[scale=#1, >=stealth, x=1.25cm, y=0.65cm]
    \draw[->] (-2.0,0) -- (2.05,0) node[right] {$x$};
    \draw[->] (0,-2.6) -- (0,2.6) node[above] {$y$};
    \draw[dashed, gray] (-2.0,2) -- (2.0,2);
    \draw[dashed, gray] (-2.0,-2) -- (2.0,-2);
    \draw[thick, domain=-1.9:1.9, samples=120]
      plot (\x,{2*sin(deg(2*\x+pi/6))});
    \node[left] at (0,2) {$2$};
    \node[left] at (0,-2) {$-2$};
    \node[below left] at (0,0) {$O$};
  \end{tikzpicture}%
}

\newcommand{\qTrigGraphPropertiesAnswer}{AD}

\newcommand{\qTrigGraphPropertiesSolution}{%
  A. 角频率为 \(\omega=2\)，所以最小正周期
  \[
    T=\frac{2\pi}{|\omega|}=\pi,
  \]
  正确。\\
  B. 正弦函数在相位满足 \(2x+\frac\pi6=\frac\pi2+k\pi\) 时取得极值，对称轴为
  \[
    x=\frac\pi6+\frac{k\pi}{2},
  \]
  因此 \(x=\frac\pi6\) 确实是其中一条对称轴，B 正确。\\
  C. 令 \(t=2x+\frac\pi6\)。当 \(x\in\left[-\frac\pi3,\frac\pi6\right]\) 时，
  \[
    t\in\left[-\frac\pi2,\frac\pi2\right].
  \]
  \(\sin t\) 在该区间单调递增，所以 C 正确。\\
  D. 将 \(f(x)\) 的图像向右平移 \(\frac\pi{12}\) 个单位，相当于把 \(x\) 替换为 \(x-\frac\pi{12}\)：
  \[
    2\sin\left(2x-\frac\pi6+\frac\pi6\right)=2\sin2x,
  \]
  所以 D 正确。\\
  按当前题面四项均成立，正式作为多选题使用前应核对原始选项设计。%
}

\newcommand{\qTrigGraphPropertiesDiagnosis}{%
  三角函数图像题要把“周期、对称轴、单调区间、平移”分别翻译成相位条件，不能只靠图像印象。%
}

\newcommand{\qTrigGraphPropertiesTeachingNotes}{%
  适合训练统一相位变量 \(t=2x+\frac\pi6\) 的方法。通过一个相位变量同时处理周期、对称、单调与平移，可显著减少符号混乱。%
}
